\section{Protezione dei dati clinici}\label{sec:dati}

Questa sezione riguarda i dati che restano fuori dalla catena: i referti, il registro dettagliato degli accessi e le chiavi che servono a leggerli. L'obiettivo è che chi legge il database, per esempio da un backup sottratto, non ottenga alcun dato sanitario.

\subsection{Cifratura dei referti}
\paragraph{Schema a busta.} Ogni referto è cifrato con AES-256-GCM, che garantisce riservatezza e integrità\cite{sp800-38d}, usando una propria chiave casuale di 256 bit (DEK). La DEK è cifrata a sua volta con una chiave di cifratura delle chiavi (KEK) ricavata dalle chiavi dell'archivio (\emph{envelope encryption}): ogni chiave cifra un solo documento, e per cambiare le chiavi dell'archivio basta cifrare di nuovo le DEK. Si cifra anche il titolo, perché un titolo come ``Esame HIV'' è già un dato sanitario; in chiaro restano solo identificativo, paziente e data.

\paragraph{Cifratura ibrida post-quantistica.} I referti vanno protetti per decenni, e un avversario potrebbe conservarli oggi per decifrarli quando avrà un computer quantistico (\emph{harvest now, decrypt later}). La KEK nasce quindi da due meccanismi di incapsulamento della chiave (KEM) usati insieme:
\begin{itemize}
  \item \textbf{ML-KEM-768}, lo standard post-quantistico del NIST basato su reticoli (FIPS~203\cite{fips203});
  \item \textbf{X25519}, nella forma DHKEM di RFC~9180\cite{rfc9180}, un meccanismo classico su curve ellittiche.
\end{itemize}
Per leggere un referto bisogna rompere entrambi: se si scoprisse una debolezza in ML-KEM, ancora giovane, resterebbe X25519; contro un computer quantistico resterebbe ML-KEM. I due segreti sono combinati con HKDF-SHA256 secondo la raccomandazione NIST SP~800-227\cite{sp800-227}, che chiede di includere nella derivazione anche i testi cifrati e le chiavi pubbliche del destinatario (listato~\ref{lst:kek}). Lo schema ha un numero di versione, incluso nella derivazione, che permette di cambiarlo in futuro senza rendere illeggibili i documenti esistenti; chi alterasse la versione dichiarata in un documento otterrebbe una KEK diversa e la decifratura fallirebbe.

\begin{lstlisting}[language=TypeScript,caption={Derivazione della chiave di cifratura delle chiavi (\texttt{crypto/vault.ts}, estratto).},label={lst:kek}]
const ikm = Buffer.concat([mlkemSecret, ecSecret]);
const transcript = createHash("sha256").update(kemCt).update(ecCt);
transcript.update(recipient.mlkem.export({ format: "der", type: "spki" }));
transcript.update(recipient.x25519.export({ format: "der", type: "spki" }));
const info = Buffer.concat([Buffer.from(`cartella-clinica/kek/v${version}`), transcript.digest()]);
return Buffer.from(hkdfSync("sha256", ikm, Buffer.alloc(0), info, 32));
\end{lstlisting}

\paragraph{Legame con il contesto.} Come dati associati di AES-GCM (AAD) si usano identificativo del documento e paziente: un referto copiato nella cartella di un altro paziente, o due testi cifrati scambiati, non si decifrano, e il backend consegna gli altri referti segnando quello come \emph{integrità non verificata}. I referti si decifrano solo in memoria, solo per un medico autorizzato e solo dopo che l'accesso è stato scritto sulla catena. Tutti gli algoritmi sono quelli di \texttt{node:crypto} (OpenSSL): nessun algoritmo crittografico è scritto in proprio. DEK e KEK vengono sovrascritte con zeri appena non servono più.

\subsection{Firma dei referti}\label{sec:firma}
La cifratura non dice chi ha scritto un referto. Il medico lo firma con la stessa passkey con cui accede: invece di una challenge casuale, il server chiede all'autenticatore di firmare l'impronta SHA-256 del contenuto (paziente, autore, titolo, testo, data e un valore casuale, preceduti da un'etichetta che distingue questo uso da ogni altro, listato~\ref{lst:firma}). Firmando la challenge dopo impronta o PIN, la passkey firma il referto; se la passkey usa ML-DSA, la firma è post-quantistica.

\begin{lstlisting}[language=TypeScript,caption={Impronta del contenuto da firmare (\texttt{crypto/signature.ts}, estratto).},label={lst:firma}]
const canonical = JSON.stringify(["cartella-clinica/referto/v1",
  p.patient, p.author, p.title, p.body, p.signedAt, p.nonce]);
return Uint8Array.from(createHash("sha256").update(canonical).digest());
\end{lstlisting}

Il server prepara la firma solo per un medico con un consenso valido \emph{per la cura} del paziente. La richiesta di firma vale cinque minuti e può essere consumata una sola volta e solo dal medico che l'ha aperta; challenge di login e richieste di firma sono conservate separatamente, così una risposta di login non vale come firma. Come per la lettura, la scrittura passa dal registro di audit, e solo se la registrazione riesce il referto viene cifrato e salvato con la prova della firma, compresa la chiave pubblica usata, così la firma resta verificabile anche dopo la revoca della passkey. A ogni consegna il backend verifica di nuovo la firma: \emph{valida}, \emph{non valida} (il contenuto non è più quello firmato) o \emph{assente} (referti caricati dal laboratorio).

La firma lega il contenuto alla passkey e, attraverso la registrazione, al medico: chi ottenesse le chiavi dell'archivio potrebbe riscrivere un referto, ma non firmarlo. Il legame tra chiave e persona è però garantito dal nostro database e non da un certificato di un ente terzo, quindi la firma non ha il valore legale di una firma digitale qualificata.

\subsection{Registro degli accessi cifrato}
Sulla catena ogni accesso compare solo come impronta. I dettagli (chi, per quale paziente, per quale finalità, con quale esito e, in emergenza, la motivazione) sono dati personali e clinici, conservati nel database solo cifrati con le chiavi dell'archivio; come dato associato si usa l'impronta, cioè il valore scritto sulla catena, e alla lettura l'impronta viene ricalcolata per scoprire alterazioni. Il registro si consulta dalla pagina del paziente (sezione~\ref{sec:registro-paziente}) o, per il responsabile della protezione dei dati, con lo script \texttt{npm run registro-accessi}. Contiene anche i tentativi d'emergenza respinti, che sulla catena non lasciano traccia.

\paragraph{Cercare senza rivelare.} Per trovare le voci di un paziente senza decifrare l'intero registro, ogni voce porta un'\emph{etichetta}: l'HMAC-SHA256 del nome utente, con una chiave derivata con HKDF da quelle dell'archivio. Chi legge il database non può ricondurre le etichette a un nome, nemmeno provando tutti i nomi utente, perché senza la chiave non può calcolarle. L'etichetta serve solo a cercare: fa fede il paziente scritto nei dettagli cifrati, così una voce spostata sotto un'altra etichetta viene mostrata solo come alterata, senza rivelarne il contenuto. Restano visibili il numero di voci per etichetta e le loro date.

\subsection{Cancellazione su richiesta}\label{sec:cancellazione}
Il diritto alla cancellazione (art.~17 GDPR\cite{gdpr}) ha eccezioni, tra cui gli obblighi di legge e la difesa di un diritto: la documentazione clinica ha obblighi di conservazione lunghi, e il registro degli accessi serve a rispondere di un abuso. La procedura, eseguita dalla struttura con \texttt{npm run cancella-utente} dopo aver identificato chi la chiede e con una conferma esplicita, distingue quindi tre gruppi di dati:
\begin{itemize}
  \item \textbf{identità digitale e profilo}, sempre cancellati: passkey, sessioni, inviti, richieste in sospeso e il profilo con l'indirizzo sulla catena;
  \item \textbf{referti}, conservati cifrati, salvo indicazione esplicita quando la conservazione non è più dovuta;
  \item \textbf{registro degli accessi}, conservato cifrato, per l'accountability.
\end{itemize}
Il punto centrale è cancellare il legame tra la persona e il suo indirizzo sulla catena, come indicano le linee guida europee\cite{edpb-blockchain}: i dati sulla catena non si cancellano, ma sono impronte e uno pseudonimo che, senza quel legame, non riconducono più a nessuno. Senza profilo nessun medico può più aprire la cartella. Uno username cancellato non può essere riassegnato, altrimenti una persona nuova erediterebbe i referti conservati di un'altra: la procedura conserva l'impronta SHA-256 del nome per bloccarne il riuso, una pseudonimizzazione e non un anonimato, perché un nome si può indovinare. Sette test, di cui uno su MongoDB, verificano tutti questi comportamenti.

\subsection{Custodia delle chiavi}\label{sec:chiavi}
Chi possiede le chiavi dell'archivio legge tutto, e chi le perde perde tutto. La gestione segue le indicazioni del NIST sulla gestione delle chiavi\cite{sp800-57} con due misure.

\paragraph{Chiavi cifrate con una passphrase.} Il file \texttt{secrets/archivio.key}, escluso da Git e leggibile solo dal proprietario, contiene le chiavi private cifrate con AES-256-GCM, con una chiave ricavata dalla passphrase tramite Argon2id (RFC~9106\cite{rfc9106}: 64~MiB di memoria, 3 passate). Argon2id è \emph{memory-hard}: ogni tentativo di indovinare la passphrase costa tempo e memoria anche su hardware specializzato. La passphrase, di almeno 12 caratteri, non è salvata da nessuna parte e viene chiesta all'avvio; i parametri letti dal file sono accettati solo entro limiti, così un file manipolato non può né indebolirli né bloccare l'avvio. Argon2id non è tra le funzioni approvate dal NIST, che per derivare chiavi da una password prevede PBKDF2 (SP~800-132\cite{sp800-132}); è stato scelto, come raccomandano RFC~9106 e OWASP\cite{owasp-password}, perché PBKDF2 non richiede memoria e si accelera molto di più con le GPU. Dove serve la conformità FIPS andrebbe usato PBKDF2 con molte iterazioni.

\paragraph{Backup a soglia.} Un backup in un solo posto è anche una sola copia da rubare. Il backup usa lo schema a soglia di Shamir\cite{shamir}: le chiavi sono cifrate con una chiave casuale, divisa in 5 quote, di cui 3 qualsiasi bastano a ricostruirla, mentre 2 non rivelano nulla, per una proprietà matematica e non per mancanza di potenza di calcolo. Ogni quota va a un custode diverso (direzione sanitaria, responsabile della protezione dei dati, sistemi informativi, cassaforte, sede di disaster recovery). L'algoritmo viene da una libreria verificata da due audit indipendenti\cite{privy-sss}, che però non riconosce quote errate: per questo si divide una chiave casuale e non le chiavi stesse, così una quota alterata produce una chiave sbagliata che AES-GCM rifiuta. Lo script di backup rifiuta di creare le quote dentro il progetto, e Git le ignora ovunque si trovino.

\subsection{Verifica}
I test su cifratura e custodia delle chiavi sono 34. Verificano, tra l'altro, che il testo in chiaro non compaia nel documento cifrato, che l'alterazione di ciascun campo cifrato venga scoperta, che un documento spostato su un altro paziente non si decifri e che basti \emph{una sola} delle due chiavi private sbagliata perché la decifratura fallisca, prova che lo schema usa davvero entrambe le componenti; per la custodia, passphrase giusta e sbagliata, file alterato o con parametri fuori limite, ricostruzione con qualunque terna di quote e rifiuto con due quote o con una quota alterata. I 13 test sulla firma dei referti, di cui uno su MongoDB, verificano firme classiche e ML-DSA e, tra gli attacchi, la scrittura senza consenso, la firma con la passkey di un altro, il riuso di una firma intercettata, una risposta di login presentata come firma e la modifica del testo dopo la firma.

\subsection{Rischi residui}
\begin{itemize}
  \item \textbf{Chiavi in memoria.} A server acceso le chiavi sono nella memoria del processo: le misure descritte proteggono file e backup, non la macchina accesa. La soluzione è un modulo hardware di sicurezza (HSM) o un servizio di gestione delle chiavi.
  \item \textbf{Una sola chiave per l'archivio.} Chiavi separate per paziente limiterebbero l'impatto di una compromissione e permetterebbero di rendere illeggibili i referti di una persona cancellandone solo la chiave (\emph{crypto-shredding}), anche nei backup.
  \item \textbf{Metadati e spazio.} Paziente e data dei referti restano in chiaro per poterli ritrovare; i campi cifrati, salvati come testo base64, occupano un terzo in più del necessario (sezione~\ref{sec:valutazione}).
\end{itemize}
